\section*{Capítulo \ref{ch:b2:fragment-orbitals} --- Orbitais de fragmentos e moléculas poliatômicas}

\begin{solution}{exo:b2:fragment-orbitals:1}
$h_1 + h_2$ e $h_1 - h_2$. O plano $xz$, que contém a bissetriz do ângulo e é
perpendicular à molécula, troca os dois hidrogênios: $h_1 - h_2$ muda de
sinal, $h_1 + h_2$ não muda.
\end{solution}

\begin{solution}{exo:b2:fragment-orbitals:2}
Seis orbitais de valência (o 2s e os três 2p do oxigênio, os dois 1s dos
hidrogênios), oito elétrons de valência, quatro orbitais preenchidos.
\end{solution}

\begin{solution}{exo:b2:fragment-orbitals:3}
Os níveis de valência preenchidos são um orbital $a_1$ e um conjunto de três
orbitais $t_2$ degenerados: duas energias, duas bandas. A banda $t_2$ retira
elétrons de três orbitais que contêm seis elétrons, contra um orbital e dois
elétrons para $a_1$.
\end{solution}

\begin{solution}{exo:b2:fragment-orbitals:4}
A ligação $\pi$ exige que os dois orbitais p dos carbonos sejam paralelos, o
que coloca os dois planos \ce{CH2} em um plano comum.
\end{solution}

\begin{solution}{exo:b2:fragment-orbitals:5}
A amônia perde um elétron de seu par não ligante, um orbital em grande parte
não ligante centrado no nitrogênio; o metano, de um \omterm{def:b2:diatomic-mos:bonding}{orbital ligante} $t_2$. Um
elétron não ligante fica mais alto, mais perto do nível atômico, que um
elétron ligante: ele é retirado mais facilmente.
\end{solution}

\begin{solution}{exo:b2:fragment-orbitals:6}
A água linear tem dois orbitais p não ligantes preenchidos e degenerados; a
flexão abaixa um deles e eleva ligeiramente um nível ligante: um ganho líquido
com oito elétrons. \ce{BeH2} tem apenas quatro, que preenchem os dois \omterm{def:b2:diatomic-mos:bonding}{orbitais
ligantes}; o orbital que a flexão abaixaria está vazio, enquanto os níveis
ligantes perdem sobreposição: a forma linear é a melhor.
\end{solution}

\begin{solution}{exo:b2:fragment-orbitals:7}
Cada ligação C--H ou C--C em um carbono vizinho do centro catiônico pode se
alinhar com o orbital p vazio e dá cerca de $2\beta^2/\Delta$ de estabilização.
\ce{CH3+} não tem tal vizinho; os cátions etila, isopropila e terc-butila têm
um, dois e três grupos metila junto ao centro, cada um oferecendo uma ligação
alinhada: a estabilidade cresce nessa ordem.
\end{solution}

\begin{solution}{exo:b2:fragment-orbitals:8}
Para dois níveis iguais, o desdobramento é $2|\beta|$, e os dois elétrons $\pi$
ganham $2|\beta|$, proporcional a $\cos\theta$. Esse ganho cai à metade em
$\theta = 60^\circ$ e se anula em $90^\circ$.
\end{solution}

\begin{solution}{exo:b2:fragment-orbitals:9}
O borano tem seis elétrons de valência: eles preenchem os três \omterm{def:b2:diatomic-mos:bonding}{orbitais
ligantes}, e o orbital p do boro perpendicular ao plano fica vazio; nada se
ganha com a piramidalização. A amônia tem dois a mais: eles ocupam o orbital
que é não ligante na forma plana e, como na água, a piramidalização permite
que ele se misture com o orbital 2s e com a combinação dos hidrogênios e desça.
\end{solution}

\begin{solution}{exo:b2:fragment-orbitals:10}
As equações seculares para $c_1 = c_2 = c_3$ dão $\alpha + 2\beta$; as duas
combinações ortogonais a ela dão $\alpha - \beta$ (degeneradas). Dois elétrons em
$\alpha + 2\beta$: energia
$2\alpha + 4\beta$, mais baixa que a de $\ce{H2} + \ce{H+}$ ($2\alpha + 2\beta$), pois $\beta < 0$.
\end{solution}

\begin{solution}{exo:b2:fragment-orbitals:11}
$\pi_1$: os três coeficientes de mesmo sinal, nenhum nó (ligante); $\pi_2$: sinais
opostos nos carbonos das extremidades e um nó no centro (não ligante); $\pi_3$:
sinais alternados, dois nós (antiligante). O ânion tem quatro elétrons $\pi$:
$\pi_2$ é seu nível ocupado mais alto, com densidade apenas nos carbonos das
extremidades.
\end{solution}

\begin{solution}{exo:b2:fragment-orbitals:12}
Em cada um dos dois planos que contêm o eixo, três orbitais p paralelos (O, C, O)
dão um \omterm{def:b2:diatomic-mos:bonding}{orbital ligante}, um não ligante (apenas nos oxigênios) e um antiligante:
seis orbitais $\pi$ ao todo. Os dois ligantes e os dois não ligantes estão
preenchidos: oito elétrons $\pi$ (duas ligações $\pi$ e dois pares não
ligantes do oxigênio na representação de Lewis).
\end{solution}

\begin{solution}{pb:b2:fragment-orbitals:1}
\textbf{1.} Oito: seis do oxigênio, um de cada hidrogênio.
\textbf{2.} $h_1 + h_2$ e $h_1 - h_2$.
\textbf{3.} 2s e $2\mathrm p_z$.
\textbf{4.} $2\mathrm p_y$.
\textbf{5.} $2\mathrm p_x$, perpendicular ao plano molecular.
\textbf{6.} Seis; quatro preenchidos.
\textbf{7.} $h_1 - h_2$, que muda de sinal, como $2\mathrm p_z$, através do plano
perpendicular ao eixo que passa pelo oxigênio.
\textbf{8.} $h_1 + h_2$.
\textbf{9.} $2\mathrm p_x$ e $2\mathrm p_y$, perpendiculares ao eixo: um par não ligante
degenerado.
\textbf{10.} (2s, $h_1 + h_2$) ligante$^2$, ($2\mathrm p_z$, $h_1 - h_2$) ligante$^2$, depois
$2\mathrm p_x^2\,2\mathrm p_y^2$.
\textbf{11.} Não: o par degenerado está cheio, todos os elétrons estão emparelhados.
\textbf{12.} Na energia de um orbital 2p do oxigênio: ele é não ligante.
\textbf{13.} $2\mathrm p_y$ (com o eixo da molécula linear ao longo de $z$ e a flexão em $yz$).
\textbf{14.} Ele se mistura com o 2s e com $h_1 + h_2$ e desce: o nível $3a_1$.
\textbf{15.} O nível ligante ($2\mathrm p_z$, $h_1 - h_2$), que perde um pouco de
sobreposição quando os hidrogênios saem do eixo: ele se torna $1b_2$.
\textbf{16.} O nível preenchido que desce ganha mais do que o outro perde: com
oito elétrons, a forma angular é a mais estável.
\textbf{17.} $104.48^\circ$, entre $90^\circ$ (ligações formadas apenas por orbitais p) e
$109.5^\circ$: o orbital 2s participa em parte da ligação.
\textbf{18.} $2\mathrm p_x$, que continua sendo o não ligante $1b_1$.
\textbf{19.} Quatro.
\textbf{20.} Do $1b_1$, o orbital $2\mathrm p_x$ não ligante do oxigênio.
\textbf{21.} Retirar um elétron não ligante quase não muda as ligações: o íon
tem a geometria da molécula, e pouca energia vibracional é excitada. Retirar
um elétron ligante muda os comprimentos e os ângulos de ligação, e a banda se
espalha por muitos níveis vibracionais.
\textbf{22.} \qty{12.62}{eV} contra \qty{13.62}{eV}: o oxigênio da água tem uma
carga parcial negativa tomada dos hidrogênios, o que eleva seus níveis 2p.
\textbf{23.} Eles não são dois níveis iguais: suas duas combinações são os
orbitais $3a_1$ e $1b_1$, de energias diferentes. As duas imagens descrevem a
mesma densidade eletrônica total; só os orbitais deslocalizados dão as
energias de ionização.
\textbf{24.} A água tem \textbf{quatro} bandas de ionização de valência, a
mais baixa a $\boldsymbol{\approx \qty{12.6}{eV}}$.
\end{solution}
